\section{Exact overlap corrections for rectified sums}
\label{loc:section}

The first constructive case turns the overlapping-shadow idea into an exact
local record. A merge can preserve a common component while removing a
contrast that still matters to the target. The proposal is to identify that
missing contribution before the merge and make it available where the
adjacent computation needs it. This formalizes the selective interpretation
developed in the discussion~\cite{discussion}. The results below concern
specified functions and real arithmetic; they do not assume that a learned
selector already knows which features to preserve.

\subsection{Projection loss and rectification loss are different}
\label{loc:two_losses}

Write $\rho(t)=\ReLU(t)=\max(t,0)$. A linear projection $P$ loses every
distinction in its kernel: $P(x+v)=Px$ whenever $Pv=0$. Rectification can then
lose further distinctions because all nonpositive preactivations produce
zero. These are separate operations, with separate requirements for a repair.

The CReLU construction retains the positive and negative parts of a
measurement~\cite{crelu}:
\begin{equation}
 \CReLU(z)=
 \begin{pmatrix}\rho(z)\\\rho(-z)\end{pmatrix},
 \qquad z=\rho(z)-\rho(-z).
 \label{loc:crelu_inverse}
\end{equation}
It therefore preserves the supplied measurement $z$. If $z=Px$, the identity
does not recover a component of $x$ already removed by $P$.

In the motivating pair collapse, let
\begin{equation}
 S=A+B,\qquad D=A-B,\qquad
 H=\rho(S),\qquad F=\rho(A)+\rho(B).
 \label{loc:pair_definition}
\end{equation}
Here $D$ is a contrast lost by the sum. The inputs $(1,-1)$ and $(0,0)$ both
have $S=0$, so even $\CReLU(S)$ gives the same record. Their required targets
are nevertheless $1$ and $0$. This is the distinction identified in the
earlier projection example~\cite{projection}. A record intended to preserve
it must branch before the sum removes it. Merely branching before the
subsequent ReLU is too late.

\subsection{The cancellation correction}
\label{loc:pair_correction}

\begin{proposition}[Exact local correction]
\label{loc:pair_prop}
For every $A,B\in\R$, define
\begin{equation}
 \kappa(A,B)=\rho(A)+\rho(B)-\rho(A+B).
 \label{loc:kappa_definition}
\end{equation}
Then $F=H+\kappa(A,B)$, and the following formulas are equivalent:
\begin{align}
 \kappa(A,B)
 &=\frac{|A|+|B|-|A+B|}{2}
 \label{loc:kappa_absolute}\\
 &=\frac12\rho\bigl(|D|-|S|\bigr)
 \label{loc:kappa_contrast}\\
 &=\begin{cases}
       0,&AB\ge0,\\
       \min(|A|,|B|),&AB<0.
    \end{cases}
 \label{loc:kappa_gate}
\end{align}
In particular, $\kappa$ is nonnegative.
\end{proposition}

\begin{proof}
Use $\rho(t)=(t+|t|)/2$ in \eqref{loc:kappa_definition}. The linear terms
cancel, giving \eqref{loc:kappa_absolute}; the triangle inequality gives
nonnegativity. If $A,B$ have the same sign, $|S|=|A|+|B|\ge|D|$. If their
signs differ, $|D|=|A|+|B|\ge|S|$. Hence
\[
 |A|+|B|=\max(|S|,|D|),
\]
which yields \eqref{loc:kappa_contrast}. In the opposite-sign case,
$|S|=\bigl||A|-|B|\bigr|$, so
\[
 \kappa=\frac{|A|+|B|-\bigl||A|-|B|\bigr|}{2}
        =\min(|A|,|B|).
\]
The same-sign case gives zero, including either input being zero.
\end{proof}

The gate identifies precisely when the collapsed sum loses a contribution
needed by this target. A binary gate alone is insufficient: the missing
amount varies. The record $\kappa$ supplies that amount. Conditional on the
retained value $H$, the transformations $F=H+\kappa$ and $\kappa=F-H$ are
inverse. Thus this record contains exactly the missing target value
conditional on $H$, without specifying which input supplied the positive
contribution or reconstructing the original pair.

Keeping a signed contrast alongside a rectified sum is a different proposal,
and does not suffice in general. The following collision has strictly
negative sums:
\begin{center}
\begin{tabular}{rrrrrr}
\toprule
$A$&$B$&$S$&$H$&$D$&$F$\\
\midrule
$1$&$-4$&$-3$&$0$&$5$&$1$\\
$2$&$-3$&$-1$&$0$&$5$&$2$\\
\bottomrule
\end{tabular}
\end{center}
The retained pair $(H,D)$ is identical while the target differs. Retaining
$(S,D)$ would recover $A=(S+D)/2$ and $B=(S-D)/2$. For this symmetric target,
$(S,|D|)$ also suffices. The correction $(H,\kappa)$ provides a more direct
statement of what the target requires. None of these facts establishes a
computational advantage for one implementation.

The distinction between a gate and a record can also be checked on a whole
family of inputs. Set $(A,B)=(t,-t)$ for $0<t<1$. Every member has $H=0$ and
the same opposite-sign gate, but its target and correction are both $t$.
An exact record must distinguish these target values; a flag saying that
cancellation occurred cannot do so. Conversely, $(t,-t)$ and $(-t,t)$ have
opposite signed contrasts but the same target. This target does not require
the record to distinguish them. The example identifies what must vary in
the record without demanding preservation of an unnecessary distinction.

\subsection{Repairing the original four-feature teacher}
\label{loc:teacher_repair}

In the motivating QDO setup, write the four intrinsic teacher features as
$u_1,\ldots,u_4$. The teacher is~\cite{qdo}
\begin{equation}
 f_\star(x)=
 \frac{-\rho(u_1)+\rho(u_2)-\rho(u_3)+\rho(u_4)}{2}.
 \label{loc:teacher}
\end{equation}
Group the features by their output coefficient sign. Define
\begin{align}
 S_+&=u_2+u_4,&\kappa_+&=\kappa(u_2,u_4),\\
 S_-&=u_1+u_3,&\kappa_-&=\kappa(u_1,u_3).
\end{align}
Applying Proposition~\ref{loc:pair_prop} to each pair gives the exact repair
\begin{equation}
 \boxed{
 f_\star(x)=
 \underbrace{\frac{\rho(S_+)-\rho(S_-)}{2}}_{f_{\rm pair}(x)}
 +\frac{\kappa_+-\kappa_-}{2}.}
 \label{loc:teacher_repair_identity}
\end{equation}
Indeed, the positive pair contributes $\rho(S_+)+\kappa_+$ and the negative
pair contributes $-\rho(S_-)-\kappa_-$. Adding them and dividing by two
reproduces every term of \eqref{loc:teacher}. Each local record is
nonnegative; their contribution to the scalar output has the teacher's
signed coefficients.

This construction identifies sufficient information for the collapsed-pair
example without requiring recovery of all eight ambient input coordinates.
It is an oracle construction when the teacher features $u_j$ are supplied.
A student that has not learned those features cannot calculate these exact
records from an already collapsed representation. Accordingly, an oracle
repair experiment establishes sufficiency of the proposed information;
learning an adequate record remains a separate empirical question. Supplying
the teacher features to a record branch must be disclosed as additional
information in any comparison.

\subsection{Position weights and signed readouts}
\label{loc:weights}

Suppose the required pair target assigns known positive position weights:
\[
 F_\alpha=\alpha_0\rho(A)+\alpha_1\rho(B),
 \qquad 0<\alpha_0<\alpha_1.
\]
The later or upper position receives the larger weight. Positive homogeneity,
$\rho(\alpha A)=\alpha\rho(A)$ for $\alpha>0$, allows the substitutions
$u=\alpha_0A$ and $v=\alpha_1B$. Therefore
\begin{equation}
 F_\alpha=
 \rho(\alpha_0A+\alpha_1B)+\kappa(\alpha_0A,\alpha_1B).
 \label{loc:ordered_pair}
\end{equation}
The common/contrast formulas apply to these weighted inputs. The weights
belong to positions if moving a layer is intended to change its influence.
This target can distinguish some reorderings, but it does not promise to
distinguish every ordering of the original data.

Negative output coefficients require care because
$\rho(\alpha A)=\alpha\rho(A)$ generally fails for $\alpha<0$.
For arbitrary real $\alpha,\beta$, set
\[
 T=\alpha\rho(A)+\beta\rho(B),\qquad S=\alpha A+\beta B.
\]
Expansion of the same ReLU identity yields
\begin{equation}
 T=\rho(S)+R_{\alpha,\beta},\qquad
 R_{\alpha,\beta}
 =\frac{\alpha|A|+\beta|B|-|\alpha A+\beta B|}{2}.
 \label{loc:signed_pair}
\end{equation}
For example, $\alpha=1$, $\beta=-1$, $A=1$, $B=2$ gives $T=-1$, a zero
rectified base, and correction $-1$. A general signed readout therefore
cannot be represented by an unqualified nonnegative correction rule.

\subsection{Arbitrary finite piles and exact local merging}
\label{loc:multiple}

Consider a fixed finite collection $u_1,\ldots,u_K$, where positive
position weights may already have been absorbed into the $u_i$. The required
target is $F_K=\sum_i\rho(u_i)$. Choose any binary merge tree. Each leaf
retains its signed value; each internal node $v$ computes
\[
 S_v=S_{v,L}+S_{v,R},\qquad
 \kappa_v=\kappa(S_{v,L},S_{v,R}).
\]
The child sums must remain signed until this operation has used them.

\begin{theorem}[Telescoping local corrections]
\label{loc:telescoping}
For every such tree,
\begin{equation}
 F_K=\rho(S_{\rm root})+
          \sum_{v\text{ internal}}\kappa_v.
 \label{loc:telescope_equation}
\end{equation}
\end{theorem}
\begin{proof}
Expand each internal record as
$\kappa_v=\rho(S_{v,L})+\rho(S_{v,R})-\rho(S_v)$.
Every nonroot internal node appears once as a child, with a positive term,
and once as a parent, with a negative term. These terms cancel. The remaining
terms are the rectified leaves minus the rectified root, proving the claim.
For one leaf the internal sum is empty and the identity also holds.
\end{proof}

The records need not all remain separate when only this aggregate is needed.
A subtree can carry a signed sum $S$ and accumulated correction $R$. Define
\begin{equation}
 (S_L,R_L)\star(S_R,R_R)=
 \bigl(S_L+S_R,\ R_L+R_R+\kappa(S_L,S_R)\bigr).
 \label{loc:merge_operator}
\end{equation}
Leaves start at $(u_i,0)$. This gives a constant number of scalar fields;
their required numerical range can still grow with the number and weights
of the leaves.

\Needspace{11\baselineskip}
\begin{proposition}[Mass coordinates and the merge law]
\label{loc:mass_coordinates}
Let $P=\sum_i\rho(u_i)$ and $N=\sum_i\rho(-u_i)$. Then
\begin{equation}
 S=P-N,\quad R=\min(P,N),\qquad
 P=\rho(S)+R,\quad N=\rho(-S)+R.
 \label{loc:mass_inverse}
\end{equation}
The merge \eqref{loc:merge_operator} is associative and commutative, with
identity $(0,0)$.
\end{proposition}
\begin{proof}
If $P\ge N$, then $S=P-N\ge0$ and $R=N$, giving both inverse formulas.
If $P<N$, then $S<0$ and $R=P$, again giving them. These formulas establish a
bijection between $(P,N)\in\R_{\ge0}^2$ and $(S,R)\in\R\times\R_{\ge0}$.
Combining two subtrees adds their $P$ and $N$ separately. The resulting
correction is
\begin{align*}
 R_{\rm new}
 &=P_L+P_R-\rho(S_L+S_R)\\
 &=R_L+R_R+\rho(S_L)+\rho(S_R)-\rho(S_L+S_R),
\end{align*}
which is exactly \eqref{loc:merge_operator}. Thus the merge is ordinary
addition in $(P,N)$ coordinates. Associativity, commutativity, and the
identity follow from addition.
\end{proof}

These algebraic properties hold in real arithmetic. Floating-point
implementations can depend on merge order: even the positive leaves
\((10^{16},1,1)\) yield different binary64 sums under the two
parenthesizations.

One immediate consequence connects the correction directly to rectified
sign channels:
\begin{equation}
 \boxed{
 \sum_i\CReLU(u_i)
 =\CReLU\!\left(\sum_i u_i\right)
   +\begin{pmatrix}R\\R\end{pmatrix}.}
 \label{loc:crelu_mass}
\end{equation}
The accumulated positive and negative masses differ from rectifying the
final signed sum by the same cancellation amount in both channels.
The individual seam corrections depend on the merge tree, but their total
$R$ does not.

Commutativity is an explicit limit on the record. The final $(S,R)$ does not
retain the original order: $(1,-1)$ and $(-1,1)$ both yield $(0,1)$. Position
weights can make a specified ordered target respond to a permutation;
source labels and transition direction must remain available when the
subsequent computation requires them. Neither weights nor this two-field
summary guarantee arbitrary future queries. In particular, changing the
readout weights after aggregation can require distinctions no longer
present.

It is useful to distinguish the information needed for the current output
from a state that supports a specified later operation. The pair
$(H,\kappa)$ gives the current target, whereas the merge rule deliberately
carries the signed $S$. For example, the subtree states $(S,R)=(-1,0)$ and
$(-3,0)$ have the same current positive mass, zero, and the same
$(\rho(S),R)$. A later request for $\rho(S+2)$ produces $1$ and $0$,
respectively. The signed sum is necessary for that additional request.
If the only request remains the sum of individually rectified leaves, one
can instead accumulate their positive mass directly. The appropriate state
therefore depends on the operations that will actually follow the overlap.
The explicit $(S,R)$ merge is closed under the signed-sum updates defined
here; its closure does not extend automatically to a new nonlinear readout,
a new set of position weights, or an unspecified future computation.

\paragraph{Signed leaf coefficients.}
For the more general target $\sum_i a_i\rho(A_i)$, initialize
\begin{equation}
 S_i=a_iA_i,\qquad
 R_i=a_i\rho(A_i)-\rho(a_iA_i)
 =\begin{cases}0,&a_i\ge0,\\a_i|A_i|,&a_i<0.\end{cases}
 \label{loc:signed_leaf}
\end{equation}
The negative case follows by writing $a_i=-c$, $c>0$, and using
$\rho(A_i)+\rho(-A_i)=|A_i|$. At every merge,
\[
 \rho(S_L+S_R)+R_{\rm new}
 =\rho(S_L)+R_L+\rho(S_R)+R_R.
\]
Induction therefore recovers the signed target exactly at the root. Here
$R$ may be negative. Omitting the leaf corrections would instead compute
$\sum_i\rho(a_iA_i)$, a different function.

The aggregate considered here can already be computed by directly summing
the rectified leaves. The derivation supplies an exact local accounting
mechanism for a collapsed representation, not evidence of a new efficient
algorithm. Its architectural value must be tested when records are learned
and transported across the actual transitions of interest.

\subsection{Numerical and backward-pass qualifications}
\label{loc:implementation}

The local correction itself has a simple stability property. Let
$\mathcal C=\{(A,B):AB\ge0\}$, the union of the two same-sign quadrants.
Then
\begin{equation}
 \kappa(A,B)=\operatorname{dist}_{\infty}\bigl((A,B),\mathcal C\bigr).
 \label{loc:distance}
\end{equation}
For opposite signs, reaching $\mathcal C$ requires moving at least one
coordinate to zero. Moving the smaller-magnitude coordinate attains
distance $\min(|A|,|B|)$; either coordinate change is a lower bound for
reaching its respective quadrant. Points already in $\mathcal C$ have
distance zero. This proves the identity. Distance to a fixed set is
one-Lipschitz, so perturbing both inputs by at most $\varepsilon$ changes
the ideal correction by at most $\varepsilon$. This is a property of the
function, not a guarantee that every algebraic implementation evaluates it
accurately.

Equivalent real formulas need not behave identically in finite precision.
For $A=-10^{20}$ and $B=1$, the mathematical correction is $1$. In binary64
arithmetic, \eqref{loc:kappa_absolute} can evaluate to zero because the
subtractions lose the small difference. The sign-gated minimum in
\eqref{loc:kappa_gate} returns $1$ for this example. Implement its gate with
sign comparisons rather than the product $AB$, which can overflow. This
avoids a particular subtraction error; signed sums, accumulated records,
and final reconstruction still require numerical error checks.

Forward identities also need a stated backward convention. Differentiating
an exact correction through its source variables gives matching ordinary
derivatives away from activation boundaries. Explicitly, when $A$, $B$, and
$S$ are nonzero,
\begin{equation}
 \frac{\partial\kappa}{\partial A}
 =\mathbf1_{A>0}-\mathbf1_{S>0},\qquad
 \frac{\partial\kappa}{\partial B}
 =\mathbf1_{B>0}-\mathbf1_{S>0}.
 \label{loc:gradient}
\end{equation}
Adding the derivatives of $H$ cancels the common $S$ terms and recovers the
two derivatives of $F$. Thus the exact source-dependent record restores the
derivatives of this specified function as well as its values at these
points. Detaching the record, replacing
it with an approximation, or choosing a different boundary convention can
change training. For example, the exact function
$\rho(z)-\rho(-z)=z$ has derivative one, but automatic differentiation with
$\rho'(0)=0$ on both branches computes zero at $z=0$. The absolute-value
version of the pair correction can also produce a wrong derivative for the
complete target at an internal cancellation boundary. At \(A=1,B=-1\),
using \(|\cdot|'(0)=\rho'(0)=0\) gives the automatic derivative
\((1/2,-1/2)\) for \(H+\kappa\), while the smooth target \(F\) has
gradient \((1,0)\). If preservation of a specified training signal is required,
that signal and its backward convention must be included in the contract
and verified explicitly.
